\section{The RMT-Cleaned Correlation Estimator}\label{sec:rmt}

Let $X \in \mathbb{R}^{T \times n}$ be the demeaned return matrix for $n$ assets
over $T$ periods, $\hat\Sigma = X^\top X / T$ the sample covariance,
$D = \mathrm{diag}(\hat\Sigma)$ the diagonal of asset variances, and
$C = D^{-1/2}\hat\Sigma\,D^{-1/2}$ the sample correlation. Individual variances are
well estimated even in the high-dimensional regime; the noise that destabilises
the portfolio lives in the \emph{correlation} structure. Random-matrix cleaning is
therefore applied to $C$, whose Marchenko--Pastur bulk is a clean band, and the
cleaned covariance is rebuilt as $\Sigma_0 = D^{1/2} C_0 D^{1/2}$. The
minimum-variance problem is solved in the standardised coordinates $y = D^{1/2} w$
(Section~\ref{sec:algorithm}), which recovers the scalar-identity Woodbury
structure and its conditioning.

\subsection{Marchenko--Pastur spectral separation}

The Marchenko--Pastur (MP) law~\cite{marchenko1967} describes the limiting spectral
distribution of a sample correlation of standardised series when $n, T \to \infty$
with $n/T \to c \in (0,1)$. Under the null of no correlation structure the bulk of
the empirical spectral distribution is supported on $[\lambda_-, \lambda_+]$ where
\[
  \lambda_{\pm} = (1 \pm \sqrt{c})^2\,,
\]
the unit-variance form ($\sigma^2 = 1$) that a correlation matrix enforces by
construction, $\mathrm{tr}(C) = n$. In finite samples $c$ is replaced by $n/T$,
giving the empirical upper edge $\hat\lambda_+ = (1 + \sqrt{n/T})^2$.

Correlation eigenvalues above $\hat\lambda_+$ carry genuine factor structure;
those below are consistent with noise and carry no reliable directional
information. We denote by $k$ the number of signal eigenvalues,
$k = |\{j : \lambda_j(C) > \hat\lambda_+\}|$. For S\&P~500 data with $n = 494$ and
$T = 1213$, $\hat\lambda_+ \approx 2.28$ and $k = 17$; for the factor-calibrated
synthetic model of Section~\ref{sec:results} the detected $k$ grows roughly as
$n/24$ in this regime. Working with the correlation rather than the covariance
matters here: the covariance mixes the (well-estimated) variances into the
spectrum, so its bulk is not a clean MP band and the edge $\hat\lambda_+$ has no
scale-free form.

\subsection{Construction of \texorpdfstring{$C_0$}{C0} and the cleaned covariance}

Let $C = U\Lambda U^\top$ be the eigendecomposition of the sample correlation.
Partition the eigenpairs by the MP edge: $U_k \in \mathbb{R}^{n \times k}$ holds
the $k$ signal eigenvectors with $\Lambda_k = \mathrm{diag}(\lambda_1,\ldots,\lambda_k)$,
and $U_0$ the noise eigenvectors. The bulk eigenvalues are replaced by their
average,
\[
  \bar\mu
  \;=\; \frac{1}{\,n-k\,}\sum_{j > k} \lambda_j
  \;=\; \frac{n - \sum_{j \leq k}\lambda_j}{n-k}\,,
\]
which preserves the trace $\mathrm{tr}(C) = n$. The resolution of the identity
$U_kU_k^\top + U_0U_0^\top = I$ yields the ``constant residual eigenvalue'' (CRE),
or eigenvalue-clipping, cleaned correlation of Laloux et
al.~\cite{laloux1999} and Bun, Bouchaud and Potters~\cite{bun2017} (popularised
for practitioners by L\'opez de Prado~\cite{lopezdeprado2018}),
\[
  C_0
  = U_k \Lambda_k U_k^\top + \bar\mu\, U_0 U_0^\top
  = \bar\mu\, I + U_k \Delta_k U_k^\top\,,
  \qquad
  \Delta_k = \Lambda_k - \bar\mu I = \mathrm{diag}(\delta_1,\ldots,\delta_k)\,,
\]
scalar-identity-plus-low-rank, with $\delta_j = \lambda_j - \bar\mu > 0$.

The cleaned covariance is $\Sigma_0 = D^{1/2} C_0 D^{1/2}$. Because $D^{1/2}$ is
positive-diagonal, $\Sigma_0 \succ 0$, and the change of variables $y = D^{1/2} w$
turns the minimum-variance problem $\min_w w^\top \Sigma_0 w$ subject to $Bw = c$,
$w \geq 0$ into
\[
  \min_{y}\; y^\top C_0\, y
  \quad\text{subject to}\quad
  B D^{-1/2}\, y = c,\quad y \geq 0\,,
\]
exactly (Section~\ref{sec:algorithm}), since $w^\top\Sigma_0 w = y^\top C_0 y$ and
$w \geq 0 \Leftrightarrow y \geq 0$. The solve therefore runs against the
scalar-identity-plus-low-rank $C_0$, not the diagonal-plus-low-rank $\Sigma_0$.

\begin{proposition}[Spectral properties of $C_0$]\label{prop:rmt_spectrum}
The cleaned correlation satisfies:
\begin{enumerate}
  \item $C_0 \succ 0$ with eigenvalues
        $\{\lambda_1, \ldots, \lambda_k, \bar\mu, \ldots, \bar\mu\}$
        and $\mathrm{tr}(C_0) = n$.
  \item $\kappa(C_0) = \lambda_{\max}(C) / \bar\mu$.
  \item $C_0 = C - U_0(\Lambda_0 - \bar\mu I)U_0^\top$: it differs from $C$ only in
        the noise subspace.
  \item The matrix-vector product $C_0 v = \bar\mu v + U_k(\Delta_k(U_k^\top v))$
        costs $O(nk)$ and requires $O(nk)$ storage.
\end{enumerate}
\end{proposition}

\begin{proof}
(1) follows from the eigendecomposition
$C_0 = [U_k\;U_0]\,\mathrm{diag}(\Lambda_k, \bar\mu I_{n-k})\,[U_k\;U_0]^\top$,
$\delta_j > 0$ for the signal eigenpairs, and the definition of $\bar\mu$; the
trace is preserved.
(2) the extreme eigenvalues are $\lambda_1$ and $\bar\mu$.
(3) is direct algebra.
(4) two dense products with $U_k \in \mathbb{R}^{n \times k}$.
\end{proof}

\begin{remark}[Condition number]\label{rem:kappa}
Cleaning lifts the noise floor of the correlation to $\bar\mu$ while preserving the
signal, so $\kappa(C_0) = \lambda_{\max}(C)/\bar\mu \approx 313$ on S\&P~500 data,
against $\kappa(C) \approx 91{,}000$ for the sample correlation (and
$\kappa(\hat\Sigma) \approx 106{,}000$ for the sample covariance): an improvement
of roughly $290\times$. Solving in the standardised coordinates $y = D^{1/2}w$
inherits this conditioning. The rebuilt covariance $\Sigma_0$ is itself far worse
conditioned ($\kappa(\Sigma_0) \approx 1{,}400$, inflated by the spread of asset
variances, a factor of $\approx 49$ between the largest and smallest on this
sample), which is precisely why the solve is posed in $y$. The single noise floor
also keeps the $k \times k$ Woodbury correction matrix well-conditioned regardless
of $n$.
\end{remark}

\begin{remark}[Relation to optimal RMT cleaning]\label{rem:rie}
CRE clipping is the simplest rotationally invariant estimator (RIE) and is not
loss-optimal: nonlinear shrinkage~\cite{ledoit2012,ledoit2020} and the oracle RIE
of Ledoit--P\'ech\'e and Bun et al.~\cite{ledoitpeche2011,bun2017}, which reshape
\emph{every} eigenvalue, dominate it in Frobenius loss. We nevertheless use CRE
for a structural reason: clipping is the only member of the family whose cleaned
correlation is \emph{scalar-identity-plus-low-rank}, the spiked
form~\cite{johnstone2001} that admits the closed-form Woodbury solve of
Section~\ref{sec:algorithm}; full-rank RIE modifications would require a general
factorisation. Whether the statistical gap matters for long-only minimum variance
is an empirical question, and on our data it is small (Section~\ref{ssec:oos}).
\end{remark}

\begin{remark}[Externally specified factor models]\label{rem:external_factors}
The Woodbury solver applies for any choice of $k$: a practitioner with an external
factor model (a commercial risk model, or a PCA with $k$ fixed by
cross-validation) can supply $k_{\mathrm{ext}}$ eigenpairs directly, bypassing the
MP threshold. What the solve requires is not the source of $k$ but the replacement
of the remaining $n - k$ correlation eigenvalues by a single scalar $\bar\mu$.
\end{remark}

\subsection{Choice of \texorpdfstring{$k$}{k}}\label{ssec:k_selection}

The threshold $k$ is determined by the MP upper edge $\hat\lambda_+$; on the
S\&P~500 sample it is stable near the central value $k = 17$. An audit at $k - 1$,
$k$, and $k + 1$ adds negligible cost (one extra Woodbury solve per audit point)
and reveals the portfolio's sensitivity to the threshold, which is small in the
bulk of the frontier (Section~\ref{sec:ksensitivity}).

Two biases are worth noting: finite-$T$ Tracy--Widom fluctuations of the largest
noise eigenvalue at scale $O(T^{-2/3})$~\cite{johnstone2001,elkaroui2008}, and
diffuse sector structure just below the edge. Both perturb $k$ by at most an
eigenvalue or two, exactly what the $k \pm 1$ audit of
Section~\ref{sec:ksensitivity} covers.
